
\section{From the Report to the Annual Meeting of the Lutheran Free Church, 1897}

Report on the Annual Meeting of the Lutheran Free Church, 1897, p. 12 ff.

The meeting in Fargo in the month of June last year forms a significant turning point in our development and—as we dare hope—in the development of our Church. The decision of June 20 to appoint a Free Church Committee, and the justification of this decision, was the fruit of a spiritual struggle of many years in our Church, and therefore also the point of departure for new spiritual movements.

Already at that time it had become clear that it was altogether indefensible to continue on the old way, forming church body against church body and party against party without end.

There had to be a change, a complete and entire reversal in the whole tendency of the development, if the Lutheran Church among our people was not to be torn apart into a multiplicity of party organizations, whose number could be reckoned only by the many small ambitions that would try to lift themselves up at the expense of the freedom of the congregations and with disregard for brotherly love.

The friends of Augsburg could have formed a church body in the old style, and with the passionate force of party spirit sought their own cause and strengthened themselves by a truly solid organization according to all the rules of the art. The temptation was great and close at hand; I confess it. It always seems so much safer to our flesh to rely on human cleverness and ability than on God and His Spirit and His Word.

Yet this was not “the good part” from the Lord. Our Church was not served by having the number of party organizations increased and discord and unfreedom grow. For peace and freedom we have labored and suffered through many years; peace and freedom were what we at least, for our part, sought in the formation of The United Church. We did not find them; just when we thought that a good beginning had been made, it became apparent that the very organization of a church body brings with it a danger which at any moment may become a fall, namely, occasion to use the united power to rule instead of to serve. The United Church succumbed to this danger. The point now was to avoid following in the same track, and the meeting in Fargo therefore demanded the following three things of its Free Church Committee when it charged it to draw up rules for a Lutheran Free Church.

The one was that the freedom of the congregation should be fully and entirely acknowledged and preserved without any restriction.

The second was that the strength of cooperation should be strengthened, not slackened.

The third was that the number of party organizations should not be increased.

It seemed a difficult and insurmountable task. And yet it is beyond doubt that this is precisely what is required if there is to be peace and freedom in the Church. We have seen that neither peace nor freedom can be won by political compromises; nor does it help to build the fences of the church bodies higher and higher, and to shut in and shut out according to arbitrary impulses and whims.

What else was there to do than to find a way by which one simply broke down “the middle wall of partition” between the parties, never to raise it again? In the sphere of doctrine, the most essential work in this direction has been done. There can scarcely be said to be any “doctrinal controversy” any more, since the Catechism has in the main prevailed over the clerical theses. By appealing to that which is held in common, the separation in doctrine has been all but overcome. What in this sphere still remains is disagreement in such doctrinal questions as directly concern the congregation and its life and freedom.

What has taken place in the sphere of doctrine must now also take place in the sphere of constitution, or of organization. Peace and freedom must be sought by taking up the cause of the congregation in full and utter earnest.

We already have many ecclesial tasks among us that are not the concern of any one church body, but simply leap over and pass through the fences of the church bodies without even seeing them, or at least without taking account of them. Some of these tasks are carried on by organizations more or less solid, and yet they work. Here one thinks of such ecclesial or Christian undertakings as, for example, the mission to the Jews, the mission to China, lay activity, the Santal mission, the deaconess cause, and in recent years also Augsburg Seminary.

Think for a moment about how this goes on among us, and perhaps it will not be so difficult after all to understand where the way of the Free Church lies.

In reality it is not possible to mistake it. It is the congregation that must come to the fore. For its life and freedom there must be labor in all church bodies and outside all church bodies. If the congregations can thrive spiritually within the church bodies, well and good; if the church bodies can tolerate the life and freedom of the congregations, well and good. We have nothing against that; there is no war between us and the church bodies, unless the church bodies bring party coercion and all its bitter consequences into the congregations. When the congregation and its life are corrupted by church-body politics, then there is struggle, utterly uncompromising struggle. For to Christian people the congregation is more precious than father and mother, than children and kin; and there is no peace between those who love the congregation and those who do it harm, whoever they may be.

If then there is to be a Free Church which is not to become a party, there is only one way: to lay the whole weight upon the congregation. It is precisely as in doctrine, to lay the weight upon the Catechism. It is no party cause. If there is any church body that does not approve of all possible effort being put forth for the congregation and its life, what sort of church body is that? It can hardly be a Christian church body. All church bodies are simply bound to acknowledge that such labor is justified, Christian, and necessary. And when the Free Church has only this one purpose, to work for the life and freedom of the congregation, then to that extent it is unassailable from an ecclesial or Christian side. Then one can—rightly—reproach us only that we do not labor enough, or do not labor in the best way. And we must indeed rejoice if anyone will help us to do it better, and most of all if those who fault us will themselves do the necessary work so much better than we are able.

But if free, independent congregations are not to lose anything of the strength of cooperation, then spiritual fellowship is needed. Free cooperation in everything that may serve to advance the Kingdom of God at home and abroad is necessary and beneficial; if this is to take place, then we must understand one another and know one another. We must exchange thought and spirit and thus have the bonds of love tied in the freedom of the Spirit. So far as I can understand, a truly free-church work will be almost impossible under our circumstances without very many free meetings and gatherings.

Such and the like have been the considerations that have also guided your Free Church Committee in its work with the proposed rules. These are now laid before you with the hope that they may be of some use in our efforts to find the way out.

It is not a proposal for a constitution. We proceed from the fact that the congregation is the true organization. These are only rules for cooperation in the one matter with which we are concerned, namely, the spiritual life and freedom of the congregations, so that they may better answer their calling and carry out their tasks.

We have therefore also given up the representative system, which among the so-called church bodies among us has been so abused. We have proceeded from the fact that in spiritual matters, those who take an interest in a cause and sacrifice something for it are also those who best understand it and most lovingly care for it.

With that, outward power also falls away altogether. Our free-church meetings have only spiritual influence, without any kind of outward bond upon the individual congregations.

Will such a thing then be able to go forward? Yes, why not, if only the congregations will it? We have labored on this voluntary way now for some years, and we cannot complain that it does not go. The rules now proposed aim at letting voluntariness continue to govern.

But in this connection one thing is the one thing needful: faith, faith in God’s Spirit, faith in His power to awaken life and preserve life, give love, create freedom, drive to labor, unite and bind the many small forces into one spiritually coherent work.

Georg Sverdrup